Panel A shows that simple news aggregates can already improve the
time-series forecast. News count, recency-weighted sentiment, and
mean-pooled CryptoBERT reduce QLIKE by 5.74\%, 6.32\%, and 6.64\%,
respectively, although their log-MSE gains are below 0.6\%.
The current-day sentiment average performs better on both losses,
reducing QLIKE by 9.03\% and log-MSE by 2.01\%. Thus a useful
news correction does not require a learned article-decay kernel,
and a higher-dimensional aggregate is not automatically more useful
than a compact sentiment summary.

Panel B shows that sentiment combined with learned decay nearly
matches the full-dimensional embedding. Its QLIKE reduction is
10.68\%, compared with 10.57\% for the refitted CryptoBERT model;
its log-MSE reduction is 2.34\%, compared with 2.48\%.
These are small numerical differences with opposite signs across
the two losses. The sentiment probability vector therefore captures
most of the improvement in the full-dimensional reference, without
establishing that the representations are equivalent in other
samples or training runs.

The PCA representation has the lowest average losses within this
table. With 256 components, QLIKE falls by 12.55\% and log-MSE
by 3.02\% relative to TS-only, exceeding the reductions from the
768-dimensional input. Compression consequently does not sacrifice
forecast accuracy in this experiment. The result also cautions
against interpreting embedding dimension as a direct measure of
the useful information delivered to the news head.

Together, the comparisons support learning how article signals
change with age, but do not place all of the forecast benefit in
that mechanism. Deterministic aggregation already helps, while
sentiment and compressed embeddings remain competitive when learned
decay is used. Panel A and Panel B also differ in aggregation and
head structure, so their contrast does not isolate encoder quality.
The evidence favors evaluating representation and temporal weighting
jointly, with the matched constant-kernel comparison providing the
more direct test of within-window discounting.
